\section{Problem Setting}
\label{sec:problem}

We study longitudinal data in which a fixed set of entities is observed repeatedly over time. Each entity carries one target quantity whose value is recorded at successive time points, together with a collection of covariates that describe the entity or its history. The aim is to divide the entities into a small number of subgroups so that entities in the same subgroup follow similar target trajectories, and so that each subgroup can be described by a short rule stated in terms of the covariates. A description of this kind lets a domain reader see both what happened, namely the shape of the trajectory, and to whom it happened, namely the covariate profile that selects the group.

\paragraph{Entities and trajectories.} Let the data contain $N$ entities indexed by $i \in \{1, \dots, N\}$. Every entity is observed on a shared grid of $T$ time points. The target trajectory of entity $i$ is the vector $\mathbf{y}_i = (y_{i,1}, \dots, y_{i,T}) \in \mathbb{R}^T$. In the case studies an entity is a country, an air-quality monitoring site, or a household, and the target is carbon dioxide emitted per capita, weekly mortality, gross domestic product per capita, fine particulate concentration, or the share of household income spent on health.

\paragraph{Descriptors.} Every entity also carries a descriptor vector $\mathbf{x}_i \in \mathbb{R}^d$. Descriptors are quantities that are known about the entity and that a reader can act on or reason about, such as demographic attributes, categorical group memberships, and summaries of how a covariate moved over time. Descriptors never include the target itself or any quantity mechanically derived from it. This restriction is stated precisely in Section~\ref{sec:algorithm} because it is what keeps a discovered rule informative rather than circular.

\paragraph{Segmentation and description.} A segmentation assigns each entity a label $\ell_i \in \{0, 1, \dots, K\}$, where label $0$ denotes the default group that collects every entity not captured by a more specific rule. The labels are produced by a rule set $\mathcal{R}$, a small ordered collection of conjunctive conditions on the descriptors, so that the label of an entity is decided by the first rule whose condition it satisfies. Because each rule is a conjunction of simple tests, the membership of any subgroup can be read as a single sentence.

\paragraph{Objective.} A useful segmentation has to meet two demands at once. Entities placed in the same subgroup should have trajectories that resemble one another more closely than the population as a whole, and the rule set that produces the subgroups should stay short. These two demands are combined into a single cost in Section~\ref{sec:algorithm}. The first is measured by comparing within-group trajectory distances against the average distance across the whole population. The second is measured by the total length of the rules. Trading one against the other yields subgroups that are coherent in their behaviour over time and simple enough to be stated and remembered.

\paragraph{Relation to neighbouring problems.} The task sits between two established ones. Trajectory clustering groups similar time series but returns clusters that are defined only by membership, so a reader is told that a group exists without being told how to recognise its members. Subgroup discovery returns rules that describe parts of a population, yet its classical formulation targets a single scalar outcome rather than a trajectory. The problem studied here keeps the descriptive rule of subgroup discovery and the temporal target of trajectory clustering, and it requires that the two be produced together rather than in separate steps.
